
\begin{frame}
    \frametitle{Заключение}
	Таким образом, нами представлен численный метод для нахождения оптимального робастного управления для систем с аддитивными неопределённостями и динамической обратной связью. 
	\vfill
	В работе приведён численный эксперимент, подтверждающий верность разработанного метода. 
	\vfill
	Использование линейных матричных неравенств для решения задач оптимального управления может расширить границы применения робастного управления для моделей, описывающих динамику шагающих роботов.
\end{frame}


\begin{frame}[plain, noframenumbering] % последний слайд без оформления
    \begin{center}
        \Huge
        Спасибо за внимание!
    \end{center}
\end{frame}

\begin{frame} % публикации на одной странице
	% \begin{frame}[t,allowframebreaks] % публикации на нескольких страницах
		\frametitle{Основные публикации}
		%% authorscopus
		\nocite{scbib1}%
		\nocite{Mistry2010}
		\nocite{RobotConfig}%
		\nocite{SAVIN2021}%
		%% authorother
		\nocite{bib1}%
		\ifnumequal{\value{bibliosel}}{0}{
			\insertbiblioauthor
		}{
			\printbibliography%
		}
	\end{frame}
